% Ad Atticum 12.3 (ad-atticum-12-03)
% Source: https://raw.githubusercontent.com/PerseusDL/canonical-latinLit/master/data/phi0474/phi057/phi0474.phi057.perseus-lat1.xml
% Fetched by scripts/fetch_latin.py

% Dateline (Perseus): Scr. in Tusculano iii Id. Iun. a. 708 (46).

\ciceroLetterOpener{CICERO ATTICO salutem}

\ciceroSection{1} unum te puto minus blandum esse quam me et, si uterque nostrum est aliquando adversus aliquem, inter nos certe numquam sumus. audi igitur me hoc ἀγοητεύτωσ dicentem. ne vivam, mi Attice, si mihi non modo Tusculanum, ubi ceteroqui sum libenter, sed μακάρων νῆσοι tanti sunt ut sine te sim tot dies. qua re obduretur hoc triduum ut te quoque ponam in eodem πάθει ; quod ita est profecto. sed velim scire hodiene statim de auctione et quo die venias. ego me interea cum libellis; ac moleste fero Vennoni me historiam non habere. sed tamen ne nihil de re, nomen illud, quod a Caesare, tris habet condiciones, aut emptionem ab hasta (perdere malo, etsi praeter ipsam turpitudinem hoc ipsum puto esse perdere) aut delegationem a mancipe annua die (quis erit cui credam, aut quando iste Metonis annus veniet?) aut Vettieni condicione semissem. σκέψαι igitur. ac vereor ne iste iam auctionem nullam faciat sed ludis factis Ἀτύπῳ subsidio currat, ne talis vir ἀλογηθῇ . sed μελήσει . tu Atticam, quaeso, cura et ei salutem et Piliae Tulliae quoque verbis plurimam.
